\documentclass[11pt]{article}

\usepackage[T1]{fontenc}
\usepackage[utf8]{inputenc}
\usepackage[margin=20mm]{geometry}
\usepackage{graphicx}
\usepackage{booktabs}
\usepackage{hyperref}

\hypersetup{colorlinks=true,linkcolor=blue,urlcolor=blue}
\graphicspath{{figures/student/}}

\title{Camera Uncertainty: A Student-Friendly Guide}
\author{}
\date{}

\begin{document}
\maketitle

\section{Why do we need a camera model?}

\textbf{What we do.} A camera records pixels, but a drop-shape model uses
millimetres. We use a camera model to turn the drop's image outline into a
physical shape; that shape can then be used to estimate surface tension.
This overview asks how uncertainty in that conversion changes the
reconstructed outline.

\textbf{Why.} A small camera error can move or reshape the outline in
millimetres. That can affect later estimates made from the drop's shape.

\begin{figure}[htbp]
  \centering
  \includegraphics[width=\textwidth,height=0.28\textheight,keepaspectratio]{01_pixels_to_surface_tension.pdf}
  \caption{Pixels need camera geometry before they can describe a physical
  drop. The drop shape in millimetres is the input to a later surface-tension
  estimate.}
\end{figure}

\section{Two parts describe the camera setup}

\textbf{What we do.} The \emph{intrinsic} settings describe the camera
itself: lens magnification, the image centre, and lens distortion. The
\emph{extrinsic} settings describe the camera's position and tilt relative
to the ArUco marker beside the drop.

\textbf{Why.} These settings come from different measurements. A checkerboard
helps estimate the camera itself; the marker helps locate the camera for the
drop measurement.

\begin{figure}[htbp]
  \centering
  \includegraphics[width=0.84\textwidth,height=0.42\textheight,keepaspectratio]{02_camera_marker_drop_setup.pdf}
  \caption{The marker provides a reference for the camera pose near the
  drop. The camera-to-marker distance in this setup is about 40 mm; the
  drawing is not to scale.}
\end{figure}

\clearpage
\section{How do we estimate uncertainty in the camera itself?}

\textbf{What we do.} The zoom5 calibration uses 49 checkerboard images, with
117 detected corners per image. The spread of the measured image errors
sets the size and direction of small corner shifts. The saved calibration
bootstrap repeats the calibration 500 times with shifted corners; the
histogram shows the resulting range of one focal length, \(f_x\).

\textbf{Why.} It is difficult to work out by hand how small corner errors
change all camera settings. Repeating the calibration with realistic small
shifts shows how much the fitted settings can vary. These are the saved
bootstrap results; no new calibrations are run for this guide.

\begin{figure}[htbp]
  \centering
  \includegraphics[width=\textwidth,height=0.46\textheight,keepaspectratio]{03_small_corner_shifts_bootstrap.pdf}
  \caption{Small corner shifts lead to a spread in the fitted focal length.
  The checkerboard and shifts on the left are an illustration; the
  histogram on the right uses the saved 500 calibration results.}
\end{figure}

\clearpage
\section{How does the marker locate the camera?}

\textbf{What we do.} The image-processing code finds the four ArUco corners,
then fits the two edges meeting at each corner. It uses six brightness
profiles along each edge and the scatter of the fitted edge positions to
estimate that corner's uncertainty. The pose bootstrap then uses the saved
corner uncertainties to solve for the camera pose 1,000 times; the marker
side is treated as exactly 10 mm.

\textbf{Why.} If the image corners move slightly, the camera pose that
explains them can move too. The faint outlines in the figure show 30 of the
1,000 saved pose draws; the corner ellipses are enlarged 20 times so they
can be seen.

The fitting code uses edge-fit residuals to estimate corner uncertainty.
No separate saved edge-residual review or pass/fail conclusion was found, so
this guide does not claim that a separate residual check passed.

\begin{figure}[htbp]
  \centering
  \includegraphics[width=0.82\textwidth,height=0.50\textheight,keepaspectratio]{04_marker_pose_bootstrap.pdf}
  \caption{Small corner uncertainty produces a small spread in the fitted
  marker outline. Yellow crosses mark the fitted corners, magenta ellipses
  are enlarged by a factor of 20, and faint outlines show 30 saved
  pose-bootstrap draws.}
\end{figure}

\clearpage
\section{What do we assume about corner errors?}

\textbf{What we do within one corner.} The code keeps the full two-direction
uncertainty for each corner. Here the measured x--y correlations are small
but not exactly zero: they are \(+0.038,-0.017,-0.008,-0.015\) for corners
1--4. The left panel shows plotting-only shifts drawn with a fixed seed.

\textbf{Why.} A corner is where two fitted edges meet. If the edges are
tilted relative to the image's horizontal and vertical axes, an error in
their positions can move the corner in both x and y. In this image the
marker is almost upright, so the two directions are nearly unrelated.

\textbf{What we do between corners.} The code treats the four corners'
random fit errors as independent. Each corner uses its own short stretches
of the two edges next to it; the stretches from neighboring corners do not
overlap.

\textbf{Why.} When pixel noise is local and the fitted regions do not share
pixels, separate random fit errors are a reasonable approximation. This
does not cover a shared effect such as blur, uneven lighting, biased edge
placement, or ink spread that moves several corners together.

\begin{figure}[htbp]
  \centering
  \includegraphics[width=\textwidth,height=0.48\textheight,keepaspectratio]{05_corner_noise_assumptions.pdf}
  \caption{The left panel shows one corner's nearly uncorrelated x and y
  shifts. The right panel shows the separate edge stretches used by the
  corner fits, with unused gaps between neighboring fits.}
\end{figure}

\clearpage
\section{What if all corners share an error?}

\textbf{What we do.} We test three fixed shifts without running Monte Carlo:
all corners move 0.1 pixel in x, all move 0.1 pixel in y, or each moves
0.1 pixel outwards so the marker looks slightly bigger. The largest median
normal contour shift is about 0.73 micrometres, compared with about
6 micrometres from the joint camera uncertainty.

\textbf{Why.} A shared effect does not disappear just because four corners
are averaged. For example, if the marker looks bigger than it really is,
the pose solver can place the camera too close. The marker size is treated
as exactly 10 mm here, so an error in its actual printed size is not
included in these pixel-shift tests.

\begin{figure}[htbp]
  \centering
  \includegraphics[width=\textwidth,height=0.46\textheight,keepaspectratio]{06_shared_corner_errors.pdf}
  \caption{An outward shared corner error makes the marker look bigger and
  can make the solved camera distance shorter. In these three 0.1-pixel
  tests, contour shifts stay below the joint camera uncertainty.}
\end{figure}

\clearpage
\section{Why compare independent and joint camera errors?}

\textbf{What we do.} In the independent calculation, a camera-intrinsic
sample and a pose sample are chosen separately. In the joint calculation,
each sampled set of camera intrinsics is used to solve the marker pose
again. The joint draws therefore keep the link between the two kinds of
camera uncertainty.

\textbf{Why.} If one camera setting changes, the pose that explains the
same marker image can change in response. The strongest saved link is
between the image-centre setting \(c_y\) and camera translation
\(\delta t_y\), with correlation about \(-0.98\). The requested
\(f_x\)--\(t_z\) pair is positively linked too (about \(+0.65\)), but it is
not the strongest pair, so the figure uses \(c_y\)--\(\delta t_y\).

\begin{figure}[htbp]
  \centering
  \includegraphics[width=\textwidth,height=0.43\textheight,keepaspectratio]{08_strongest_joint_pair.pdf}
  \caption{Independent draws have no link between the image-centre and pose
  shifts. Joint draws show that opposite shifts tend to occur together, so
  part of their effect can cancel.}
\end{figure}

\clearpage
\section{What does that mean for the drop outline?}

\textbf{What we do.} The saved virtual outline is passed through the saved
camera draws. The joint calculation gives a median camera-related normal
uncertainty of about \(6.0\,\mu\mathrm{m}\) per contour point; treating the
intrinsics and pose as independent gives about \(28.5\,\mu\mathrm{m}\).
The figure shows 30 saved outlines from each calculation, with deviations
enlarged by the same factor of 10 so they are visible.

\textbf{Why.} The joint calculation includes compensation between camera
settings. For this virtual outline, the independent calculation is about
4.7 times wider, so it is more pessimistic.

\begin{figure}[htbp]
  \centering
  \includegraphics[width=\textwidth,height=0.38\textheight,keepaspectratio]{07_independent_vs_joint_contours.pdf}
  \caption{The independent model produces a wider spread of back-projected
  outlines than the joint model. Deviations are enlarged by a factor of 10
  in both panels.}
\end{figure}

\textbf{One more important assumption.} The drop is assumed to lie exactly
in the marker plane. The camera is about 40 mm from the marker, and a depth
offset of about 0.27 mm has the same median normal effect as the joint
camera uncertainty. This offset is systematic for a fixed mount; repeated
measurements with that mount do not average it away.

\begin{figure}[htbp]
  \centering
  \includegraphics[width=0.82\textwidth,height=0.34\textheight,keepaspectratio]{09_summary_uncertainty.pdf}
  \caption{Median normal uncertainty per outline point from the saved
  analysis. The 0.27 mm depth-offset bar is the linearized effect of that
  offset; the pixel-noise value is an assumed 0.3 pixel, not a measured
  contour-detection error.}
\end{figure}

\section*{What is still simplified?}
\begin{itemize}
  \item The outline is a virtual model contour, not an edge detected from a
  real drop photograph.
  \item The drop-edge noise is assumed to be 0.3 pixel.
  \item The marker side is treated as exactly 10 mm, and only one marker
  image is used.
  \item Shared corner errors such as blur or printing bias are not included.
\end{itemize}

\end{document}
